\chapter{Prove That the Test Detects Failure}

A passing test is meaningful only if the broken system makes it fail. We will
stop the catalog service while leaving its jail and every other component
running. This isolates one integration boundary, exercises the store's error
handling, and makes recovery a single command.

Do not stop the catalog jail through Blackship for this exercise. Catalog now
has dependents, so a dependency-aware jail shutdown can also stop higher-level
services, including the test runner. We want a service-level fault inside one
otherwise-running jail.

\section{Exercise 23A: Remove the catalog API}

Open the catalog console as \texttt{root} and stop only its rc.d service:

\begin{lstlisting}[language=sh]
service catalog onestop
service catalog onestatus
\end{lstlisting}

The status command should report that catalog is not running and may exit
nonzero. Leave the catalog jail itself running.

\section{Exercise 23B: Observe the black-box failure}

Return to the test jail and run the unchanged test:

\begin{lstlisting}[language=sh]
ruby /usr/local/tests/store_e2e_test.rb
\end{lstlisting}

Expect a failure: the test wanted HTTP 200 but received HTTP 502. The request
still crossed nginx and store. Store caught the failed catalog connection and
rendered its deliberate \texttt{Bad Gateway} response, so the failure describes
the unavailable integration rather than crashing the test process.

This is the important control experiment. A test that still passed would be
using cached, stubbed, or otherwise disconnected data instead of traversing the
live catalog boundary.

\section{Exercise 23C: Restore and retest}

In the catalog console, start the service again and verify its status:

\begin{lstlisting}[language=sh]
service catalog onestart
service catalog onestatus
\end{lstlisting}

Run the same black-box test again from the test jail:

\begin{lstlisting}[language=sh]
ruby /usr/local/tests/store_e2e_test.rb
\end{lstlisting}

Expect one run, five assertions, no failures, and no errors. No test code or
endpoint changed between the failing and passing runs; only the real catalog
service state changed.

\section{Checkpoint}

Stop after the catalog outage produces a 502 test failure and restoring catalog
makes the same test pass. The experiment now demonstrates both sides of an
end-to-end contract: the deployed system succeeds when every boundary works and
the external test detects a real dependency outage when one boundary fails.
